\chapter*{Kế hoạch thực hiện công việc nhóm}
\addcontentsline{toc}{chapter}{Kế hoạch thực hiện công việc nhóm}
% \chapter* không tự gọi \markboth -- set thủ công để Header phải hiển thị đúng.
\markboth{KẾ HOẠCH THỰC HIỆN CÔNG VIỆC NHÓM}{}
\setlength{\parskip}{6pt}

Bảng dưới đây trình bày kế hoạch thực hiện công việc của \teamnameprose{} theo đúng cấu trúc mẫu chính thức của Khoa/Bộ môn (5 cột: TT, Nội dung, Người thực hiện, Thời gian thực hiện, Mức độ hoàn thành). Chi tiết bổ sung xem thêm tại \texttt{TASK\_BOARD.md} và \texttt{docs/report/bao-cao-dinh-ky-01.md} của kho mã nguồn dự án.

\begin{table}[H]
\centering
\caption{Kế hoạch thực hiện công việc nhóm}
\label{tab:ke-hoach-nhom}
\begin{tabular}{@{}p{0.05\linewidth}p{0.32\linewidth}p{0.2\linewidth}p{0.18\linewidth}p{0.17\linewidth}@{}}
\toprule
\textbf{TT} & \textbf{Nội dung} & \textbf{Người thực hiện} & \textbf{Thời gian thực hiện} & \textbf{Mức độ hoàn thành} \\
\midrule
1 & Khởi tạo scaffold dự án, kế hoạch, quy tắc làm việc (Phase 0) & Cả nhóm & Tuần 1 (06--11/07/2026) & Hoàn thành \\
2 & Nghiên cứu OWASP LLM Top 10, related work, so sánh công cụ/framework (Phase 1) & Nguyễn Văn An, Lê Đình Nghĩa & Tuần 1 (11/07/2026) & Bản nháp đầu, đang xác minh thủ công \\
3 & Phân tích yêu cầu, thiết kế kiến trúc gateway, threat model STRIDE (Phase 2) & Cả nhóm & Tuần 1 (11/07/2026) & Đã thiết kế đầy đủ (tài liệu) \\
4 & Thiết kế bộ dữ liệu red-team tổng hợp và tiêu chí đánh giá (Phase 2.5) & Nguyễn Văn An, Lê Đình Nghĩa & Tuần 1 (11/07/2026) & Đã thiết kế đầy đủ (tài liệu) \\
5 & Chuẩn bị báo cáo định kỳ 01 (nội dung và định dạng LaTeX) & Cả nhóm & Tuần 1 (11--13/07/2026) & Đang hoàn thiện \\
6 & Triển khai MVP: Gateway, Input/RAG/Output Guard, chạy đánh giá thật (Phase 3--9) & Cả nhóm & Các kỳ báo cáo tiếp theo & Chưa bắt đầu \\
\bottomrule
\end{tabular}
\end{table}

\textbf{Mốc thời gian:} Báo cáo định kỳ 01 dự kiến nộp trong khoảng 12--13/07/2026. Phase 0--2.5 (scaffold, nghiên cứu, kiến trúc/threat model, thiết kế dữ liệu đánh giá) được hoàn thành ở dạng tài liệu trước mốc này; Phase 3 (bắt đầu viết code) chỉ khởi động sau khi các phase tài liệu được xác nhận.

\textbf{Nguyên tắc phối hợp:} Mỗi giai đoạn đều phải có tài liệu và bằng chứng (evidence) đi kèm trước khi được đánh dấu hoàn thành; không mở rộng phạm vi ngoài kế hoạch nếu chưa được phê duyệt (xem thêm quy tắc trong \texttt{AGENT\_RULES.md}).
